\section{Conservative 方法}

\subsection{CQL：Conservative Q-Learning}

CQL 的思路更直接：在 Q 函数更新中加入\textbf{惩罚项}，主动压低 Q 值。

$$\hat{Q} = \arg\min_{Q} \max_\mu \; \mathbb{E}_{(s,a,s') \sim \mathcal{D}}\left[(Q(s,a) - (r + \gamma \mathbb{E}_\pi\left[Q(s', a')\right]))^2\right] + \alpha \mathbb{E}_{s \sim \mathcal{D}, a \sim \mu(\cdot|s)}\left[Q(s,a)\right] - \alpha \mathbb{E}_{(s,a) \sim \mathcal{D}}\left[Q(s,a)\right]$$

\begin{itemize}
    \item 第一项：标准 TD 更新。
    \item 第二项：压低所有动作的 Q 值（$\mu$ 覆盖整个动作空间）。
    \item 第三项：恢复数据中动作的 Q 值（不要过度压低）。
\end{itemize}

可以证明：对足够大的 $\alpha$，学到的 Q 值在策略分布下是\textbf{保守的}（低估而非高估）。

\subsection{CQL 的实现}

当使用最大熵正则化时，最优的 $\mu(a|s) \propto \exp(Q(s,a))$，此时惩罚项变为 $\log \sum_a \exp(Q(s,a))$，无需显式构造 $\mu$。

\subsection{本章小结}

CQL 通过主动压低 Q 值来防止过估计，是另一种有效的 offline RL 方法。与 IQL 的隐式约束不同，CQL 显式地构造保守估计。

